\documentclass[11pt,a4paper]{article}
\usepackage[T1]{fontenc}
\usepackage{lmodern}
\usepackage[margin=27mm]{geometry}
\usepackage{amsmath,amssymb,amsthm,mathtools}
\usepackage{microtype}
\usepackage{booktabs,array}
\usepackage{enumitem}
\usepackage[hidelinks]{hyperref}
\setlength{\parskip}{5pt plus 1pt}
\setlength{\parindent}{0pt}
\setlength{\emergencystretch}{3em}
\setlist{itemsep=2pt,topsep=3pt}
\raggedbottom
\newtheorem{theorem}{Theorem}[section]
\newtheorem{lemma}[theorem]{Lemma}
\newtheorem{corollary}[theorem]{Corollary}
\newtheorem{proposition}[theorem]{Proposition}
\theoremstyle{definition}
\newtheorem{definition}[theorem]{Definition}
\newtheorem{example}[theorem]{Example}
\theoremstyle{remark}
\newtheorem{remark}[theorem]{Remark}
\newcommand{\Z}{\mathbb{Z}}
\newcommand{\N}{\mathbb{N}}
\newcommand{\Q}{\mathbb{Q}}
\newcommand{\F}{\mathbb{F}}
\newcommand{\ps}[1]{#1[[x]]}
\newcommand{\coeff}[2]{[x^{#1}]\,#2}
\newcommand{\hl}[1]{\textbf{#1}}
\newcommand{\para}[1]{\par\medskip\noindent\textbf{#1}\enspace}
\hypersetup{pdftitle={The iterates of a self-referential power series modulo 4: a proof of four conjectures in OEIS A396099},pdfauthor={Dev D. Goyal},pdfsubject={OEIS A396099}}
\title{The iterates of a self-referential power series modulo $4$\\[6pt]\large A proof of four conjectures in OEIS A396099}
\author{Dev D. Goyal}
\date{October 2026}
\begin{document}
\maketitle
\begin{abstract}
Let $A(x)=x+x^2+\cdots$ be the power series defined by the equation
$A(x)=x+A^{\circ2}(x)\cdot A^{\circ3}(x)$, where $A^{\circ k}$ means $A$ substituted into
itself $k$ times. Its coefficients form the sequence A396099 of the On-Line Encyclopedia of Integer Sequences
(OEIS). Paul~D.~Hanna conjectured that every $a(n)$ is odd, that from $n=3$ on the remainders of $a(n)$ on division by $4$ repeat
the block $1,3,3,1$, and two similar statements about the iterates $A(A(x))$ and $A(A(A(x)))$. We prove all four. The key fact is that, modulo $4$, the series $A$ equals the explicit fraction
$\frac{x}{1-x}+\frac{2x^4}{(1-x)^3}$, and $A(A(x))$ equals $x+2x^2$. The proof covers the whole family of
equations $F=x+F^{\circ r}\cdot F^{\circ(r+1)}$, $r\ge1$. Modulo $2r$, every iterate $F^{\circ k}$ (including those
with $k<0$, built from the inverse series) equals the simple fraction $x/(1-kx)$, plus an explicit correction
$r\,x^4/(1-x)^3$ when $r$ is even and $k$ is odd. All results have been checked by computer and proved formally in the
Lean proof assistant.
\end{abstract}
\begin{center}\begin{minipage}{0.88\linewidth}\small
\textbf{How this paper was produced.} This paper reports the result of an autonomous AI research effort.
ChatGPT (OpenAI), using the GPT-6 model with Pro reasoning and working from prompts by Dev~D.~Goyal, found the proof, wrote its Lean~4 formalization, and wrote a first draft of the paper.
Dev~D.~Goyal chose the problems and asked for checks and revisions, but did not supply proof ideas.
Claude Opus~5.5, used with Claude Code, then checked the work independently: it rebuilt the formal proofs from
source and confirmed that they rely on no unproved assumptions, recomputed the sequence from its definition and tested
every claim against the published data, and read every step of the proofs. It also rewrote this paper for clarity.
Dev~D.~Goyal has verified the results and is responsible for them. The paper has not been peer reviewed.
The section ``About this paper'' at the end gives the details.
\end{minipage}\end{center}

\section{Introduction}

\subsection{The sequence}\label{sec:sequence}
The \emph{On-Line Encyclopedia of Integer Sequences} (OEIS)~\cite{oeisdb} is a large public database of integer
sequences. Mathematicians and scientists use it to identify sequences and to record what is known about them. Many
entries also record \emph{conjectures}: patterns that the contributor saw in the data but could not prove.

This paper is about entry A396099~\cite{oeis}, contributed by Paul~D.~Hanna in June 2026.
The entry defines a power series $A(x)=\sum_{n\ge1}a(n)\,x^n$ by an equation in which $A$ is substituted into
itself. The entry reads:
\begin{quote}\small\ttfamily\raggedright G.f. satisfies A(x) = x + A\^{}2(x)*A\^{}3(x).\par\smallskip Conjecture: all terms are odd.\par\smallskip Conjecture: a(n) = [1,3,3,1] repeating (mod 4) for n > 2.\par\smallskip Conjecture: [x\^{}n] A(A(x)) == 0 (mod 4) for n > 2.\par\smallskip Conjecture: [x\^{}n] ( A(x) - x*A(A(A(x))) ) == 2 (mod 4) for n > 2.\par\smallskip\hfill\textrm{--- OEIS A396099, Paul D. Hanna, June 2026}\end{quote}

Here ``G.f.''\ stands for \emph{generating function}, the series whose coefficients are the terms of the sequence.
The notation needs care. In the entry, $A^k(x)$ does not mean the $k$-th power of $A(x)$, but the $k$-th
\emph{iterate}: $A$ substituted into itself $k$ times, so that $A^2(x)=A(A(x))$ and $A^3(x)=A(A(A(x)))$. The entry
says so explicitly in its list of formulas (``$A^n(x)$ denotes the $n$-th iteration of $A(x)$''). To avoid any
confusion we write $A^{\circ k}$ for the iterate and keep $A^k$ for the ordinary power. So the sequence is defined by
\begin{equation}\label{eq:def}
A(x)=x+A^{\circ2}(x)\cdot A^{\circ3}(x),\qquad A(x)=x+a(2)x^2+a(3)x^3+\cdots,
\end{equation}
and $a(n)$ is the coefficient of $x^n$ in $A(x)$.

At first sight it is not clear that \eqref{eq:def} defines anything, since the unknown series appears on both sides,
substituted into itself many times. It does, and the coefficients can be computed one after another. Every iterate
of $A$ starts with $x$, so the product on the right starts with $x^2$; this gives $a(2)=1$. In general the
coefficient of $x^n$ on the right only involves $a(1),\dots,a(n-1)$, so each new coefficient is determined by the
earlier ones (Lemma~\ref{lem:unique} below makes this precise). The first terms, starting at $n=1$, are
1, 1, 5, 39, 383, 4357, 55057, 754223, 11028735, \dots; they grow quickly.

\subsection{The question}
The entry makes four conjectures. Here $[x^n]\,U$ denotes the coefficient of $x^n$ in a series $U$, and
``$==$'' or ``$=$ \dots{} (mod 4)'' means that the two sides leave the same remainder on division by $4$. In words:
\begin{enumerate}[label=(\alph*)]
\item every $a(n)$ is odd;
\item from $n=3$ on, the remainders of $a(n)$ on division by $4$ repeat the block $1,3,3,1$: the remainders of
$a(3),a(4),a(5),a(6)$ are $1,3,3,1$, those of $a(7),\dots,a(10)$ are again $1,3,3,1$, and so on;
\item from $x^3$ on, every coefficient of the second iterate $A(A(x))$ is divisible by $4$;
\item from $x^3$ on, every coefficient of $A(x)-x\cdot A(A(A(x)))$ leaves remainder $2$ on division by $4$.
\end{enumerate}
The table shows the data. The indices $n\ge3$, where the patterns are claimed, are printed in bold.
\begin{center}\small\setlength{\tabcolsep}{3.2pt}
\begin{tabular}{>{\raggedright\arraybackslash}p{11em}>{\centering\arraybackslash}p{1.15em}>{\centering\arraybackslash}p{1.15em}>{\centering\arraybackslash}p{1.15em}>{\centering\arraybackslash}p{1.15em}>{\centering\arraybackslash}p{1.15em}>{\centering\arraybackslash}p{1.15em}>{\centering\arraybackslash}p{1.15em}>{\centering\arraybackslash}p{1.15em}>{\centering\arraybackslash}p{1.15em}>{\centering\arraybackslash}p{1.15em}>{\centering\arraybackslash}p{1.15em}>{\centering\arraybackslash}p{1.15em}>{\centering\arraybackslash}p{1.15em}>{\centering\arraybackslash}p{1.15em}>{\centering\arraybackslash}p{1.15em}>{\centering\arraybackslash}p{1.15em}}
\toprule
$n$ & 1 & 2 & \hl{3} & \hl{4} & \hl{5} & \hl{6} & \hl{7} & \hl{8} & \hl{9} & \hl{10} & \hl{11} & \hl{12} & \hl{13} & \hl{14} & \hl{15} & \hl{16} \\ \midrule
$a(n)\bmod4$ & 1 & 1 & \hl{1} & \hl{3} & \hl{3} & \hl{1} & \hl{1} & \hl{3} & \hl{3} & \hl{1} & \hl{1} & \hl{3} & \hl{3} & \hl{1} & \hl{1} & \hl{3} \\
$\coeff{n}{A^{\circ2}}\bmod4$ & 1 & 2 & \hl{0} & \hl{0} & \hl{0} & \hl{0} & \hl{0} & \hl{0} & \hl{0} & \hl{0} & \hl{0} & \hl{0} & \hl{0} & \hl{0} & \hl{0} & \hl{0} \\
$\coeff{n}{(A-xA^{\circ3})}\bmod4$ & 1 & 0 & \hl{2} & \hl{2} & \hl{2} & \hl{2} & \hl{2} & \hl{2} & \hl{2} & \hl{2} & \hl{2} & \hl{2} & \hl{2} & \hl{2} & \hl{2} & \hl{2} \\
\bottomrule
\end{tabular}
\\[4pt]
\begin{tabular}{>{\raggedright\arraybackslash}p{11em}>{\centering\arraybackslash}p{1.15em}>{\centering\arraybackslash}p{1.15em}>{\centering\arraybackslash}p{1.15em}>{\centering\arraybackslash}p{1.15em}>{\centering\arraybackslash}p{1.15em}>{\centering\arraybackslash}p{1.15em}>{\centering\arraybackslash}p{1.15em}>{\centering\arraybackslash}p{1.15em}>{\centering\arraybackslash}p{1.15em}>{\centering\arraybackslash}p{1.15em}>{\centering\arraybackslash}p{1.15em}>{\centering\arraybackslash}p{1.15em}>{\centering\arraybackslash}p{1.15em}>{\centering\arraybackslash}p{1.15em}>{\centering\arraybackslash}p{1.15em}>{\centering\arraybackslash}p{1.15em}}
\toprule
$n$ & \hl{17} & \hl{18} & \hl{19} & \hl{20} & \hl{21} & \hl{22} & \hl{23} & \hl{24} & \hl{25} & \hl{26} & \hl{27} & \hl{28} & \hl{29} & \hl{30} & \hl{31} & \hl{32} \\ \midrule
$a(n)\bmod4$ & \hl{3} & \hl{1} & \hl{1} & \hl{3} & \hl{3} & \hl{1} & \hl{1} & \hl{3} & \hl{3} & \hl{1} & \hl{1} & \hl{3} & \hl{3} & \hl{1} & \hl{1} & \hl{3} \\
$\coeff{n}{A^{\circ2}}\bmod4$ & \hl{0} & \hl{0} & \hl{0} & \hl{0} & \hl{0} & \hl{0} & \hl{0} & \hl{0} & \hl{0} & \hl{0} & \hl{0} & \hl{0} & \hl{0} & \hl{0} & \hl{0} & \hl{0} \\
$\coeff{n}{(A-xA^{\circ3})}\bmod4$ & \hl{2} & \hl{2} & \hl{2} & \hl{2} & \hl{2} & \hl{2} & \hl{2} & \hl{2} & \hl{2} & \hl{2} & \hl{2} & \hl{2} & \hl{2} & \hl{2} & \hl{2} & \hl{2} \\
\bottomrule
\end{tabular}
\end{center}


\subsection{What this paper proves}
All four conjectures are true (Corollary~\ref{cor:app}). The reason is that, modulo $4$, the
complicated series $A$ equals the explicit series
\[
\frac{x}{1-x}+\frac{2x^4}{(1-x)^3}=x+x^2+x^3+3x^4+7x^5+13x^6+\cdots,
\]
and its second iterate $A(A(x))$ is simply $x+2x^2$. The four patterns can be read off from these formulas.

More generally, we treat the whole family of equations
\[
F=x+F^{\circ r}\cdot F^{\circ(r+1)},\qquad r=1,2,3,\dots,
\]
of which the entry is the case $r=2$. Theorem~\ref{thm:general} gives, for every $r$, every
iterate $F^{\circ k}$ of the solution modulo $2r$, for every integer $k$, positive, zero or negative. The answer is
remarkably simple: up to a small correction when $r$ is even, $F$ behaves modulo $2r$ exactly like the fraction
$x/(1-x)$, whose $k$-th iterate is $x/(1-kx)=x+kx^2+k^2x^3+\cdots$. In particular, modulo $2r$, substituting $F$ into
itself $2r$ times gives back $x$, and no smaller number of substitutions does.

\subsection{What was known}
Hanna recorded the conjectures in June 2026, as part of a family of four entries with $r=2,3,4,5$:
A396099, A396797, A396798 and A396807. For $r=3$ (entry A396797) the conjectures were proved by Shiva~Priyan
Shanmuga~Sundaram in September 2026~\cite{sundaram}; the proof is linked from that entry. Our argument for odd $r$ compares
the solution with the simple series $x/(1-x)$ and uses the uniqueness of the solution, for every odd $r$ at
once. For even $r$, as here, the comparison with $x/(1-x)$ alone does not work; the correction term $rH$ in
Theorem~\ref{thm:general} is the extra ingredient that even $r$ needs. When we checked the entries on
24~October~2026, the statements for $r=2,4,5$ were still labelled as conjectures. We are not aware of an earlier proof
of them.

\subsection{Why it is interesting}
Most equations that define sequences in the OEIS combine \emph{known} series by sums, products and substitutions.
In \eqref{eq:def} the unknown series is substituted into \emph{itself}, several times, inside its own definition.
Such self-referential equations are hard to handle: there is no closed formula for $A$, and even computing many terms
takes real effort. It is therefore striking that the remainders have simple patterns, and Hanna's conjectures ask
for an explanation.

The explanation found here is that, modulo $2r$, the complicated series $A$ is indistinguishable from the simple
fraction $x/(1-x)$, possibly plus a correction that is easy to control. Substituting a series into itself is the
basic operation of \emph{iteration}, that is, of repeating a process, and iterates of power series appear throughout
combinatorics and in the study of dynamical systems. Their arithmetic properties are much less understood than those
of sums and products, so complete results like the one here are rare. The method is general: guess the answer modulo
a small number, check that the guess satisfies the equation, and conclude by uniqueness. The correction needed for
even $r$ uses a second general idea, a Taylor formula that is exact because the correction squares to zero.

\subsection{How to read this paper}
Section~\ref{sec:background} explains the few notions we need: power series, substitution and iteration, inverse
series, and remainders. Section~\ref{sec:results} states the results precisely, and Section~\ref{sec:proofs} proves
them; the proof of the main theorem comes first, and the consequences for this entry are derived from it at the end of
that section. Section~\ref{sec:checks} describes how the results were checked, by computer and by a formal proof in
Lean. Section~\ref{sec:discussion} discusses the scope of the results, open questions, and related work.
\section{Background}\label{sec:background}
This section collects the definitions and facts used later. Readers who know formal power series and congruences can skip it and return when needed.
\subsection{Power series as bookkeeping}
A \emph{formal power series} is a way of storing an infinite sequence of numbers $u_0,u_1,u_2,\dots$ in one object,
\[
U(x)=u_0+u_1x+u_2x^2+u_3x^3+\cdots .
\]
The letter $x$ is only a placeholder that keeps the terms apart. We never substitute a number for $x$, so the question
of whether the series ``converges'' never arises; a power series is simply its list of coefficients. We write
$\coeff{n}{U}=u_n$ for the coefficient of $x^n$ in $U$. Two series are equal when all their coefficients are equal.

Series are added coefficient by coefficient and multiplied like polynomials:
\[
\coeff{n}{(UV)}=u_0v_n+u_1v_{n-1}+\cdots+u_nv_0 .
\]
Each coefficient of a product is a finite sum, so everything is well defined. For example,
$(1-x)(1+x+x^2+x^3+\cdots)=1$, because all the other terms cancel. So $1+x+x^2+\cdots$ is the
\emph{multiplicative inverse} of $1-x$, and we write it as $\frac{1}{1-x}$. In the same way, every series $D$ whose
constant term is $1$ has exactly one multiplicative inverse $\frac1D$. Its coefficients can be computed one at a time
from $D\cdot\frac1D=1$, using only additions, subtractions and multiplications. Every fraction in this paper has a
denominator of this kind, so it stands for an honest power series.

Two more pieces of notation. We write $U=O(x^d)$ when all coefficients of $U$ below degree $d$ are zero. So
$U-V=O(x^d)$ means that $U$ and $V$ agree up to degree $d-1$. A series of the form $x+u_2x^2+u_3x^3+\cdots$, with
constant term $0$ and coefficient of $x$ equal to $1$, is called \emph{normalized}.

\subsection{Substituting one series into another}
If $V(x)$ has constant term $0$, we can substitute $V$ for $x$ in $U$:
\[
U(V(x))=u_0+u_1V(x)+u_2V(x)^2+u_3V(x)^3+\cdots .
\]
This is called \emph{composition}, also written $U\circ V$. It makes sense because $V(x)^k$ has no terms below $x^k$,
so only the finitely many terms with $k\le n$ affect the coefficient of $x^n$. For example, if $U=x+x^2$ and $V=x^2$, then
$U(V(x))=x^2+x^4$. Composition is associative: $U\circ(V\circ W)=(U\circ V)\circ W$.

Substituting a series into itself repeatedly gives its \emph{iterates}: $U^{\circ2}(x)=U(U(x))$,
$U^{\circ3}(x)=U(U(U(x)))$, and so on, with $U^{\circ0}(x)=x$ and $U^{\circ1}=U$. They must not be confused with
\emph{powers} $U^2=U\cdot U$, $U^3=U\cdot U\cdot U$. (The OEIS sometimes writes $A^k(x)$ for the $k$-th iterate. In
this paper the small circle always means an iterate, and a plain exponent always means a power.)

\subsection{Undoing a substitution}
A normalized series can always be ``undone'' by another normalized series.

\begin{lemma}[Compositional inverse]\label{lem:inverse}
For every normalized series $U$ there is exactly one normalized series $V$ with $U(V(x))=x$, and this $V$ also
satisfies $V(U(x))=x$. If $U$ has integer coefficients, so does $V$.
\end{lemma}

We call $V$ the \emph{compositional inverse} of $U$. The OEIS calls it the \emph{series reversion} of $U$. For example,
the compositional inverse of $x-x^2$ is $C(x)=x+x^2+2x^3+5x^4+14x^5+\cdots$, whose coefficients are the Catalan numbers.
Indeed $C-C^2=x$, which says exactly that substituting $C$ into $x-x^2$ gives $x$.

\begin{proof}
Write $V=x+v_2x^2+v_3x^3+\cdots$ with unknown coefficients. In
$U(V)=V+u_2V^2+u_3V^3+\cdots$, the coefficient of $x^n$ (for $n\ge2$) equals $v_n$ plus terms that involve only
$v_2,\dots,v_{n-1}$. Setting it equal to $0$ therefore determines $v_n$ from the earlier coefficients, using only
additions, subtractions and multiplications. So there is exactly one $V$ with $U\circ V=x$, and it has integer
coefficients if $U$ does. Applying the same construction to $V$ gives a normalized $W$ with $V\circ W=x$. Then
$U=U\circ(V\circ W)=(U\circ V)\circ W=W$, so $V\circ U=V\circ W=x$ as well.
\end{proof}

\subsection{Iterates with negative index, and the order of a series}
Let $F$ be a normalized series and let $F^{-1}$ be its compositional inverse (Lemma~\ref{lem:inverse}). For a
negative integer $k$ we define $F^{\circ k}$ as the $|k|$-th iterate of $F^{-1}$; for example
$F^{\circ(-2)}=F^{-1}(F^{-1}(x))$. With this convention $F^{\circ j}\circ F^{\circ k}=F^{\circ(j+k)}$ for all integers
$j$ and $k$, just as for powers of a number. The \emph{order} of $F$ under composition is the smallest $k\ge1$ with
$F^{\circ k}(x)=x$, if there is one. Over the integers only $F=x$ has finite order, but after reduction modulo $m$
other series can have finite order, and this will happen below.

Two more pieces of notation. For a statement $S$ we write $[S]$ for $1$ if $S$ is true and $0$ if it is false
(for example $[k\text{ odd}]$). And $\binom{j}{2}=j(j-1)/2$ is the number of ways to choose $2$ objects out of $j$.

\subsection{Remainders and congruences}
For an integer $m\ge2$ we write $a\equiv b\pmod m$ (``$a$ is congruent to $b$ modulo $m$'') when $a-b$ is divisible by
$m$, that is, when $a$ and $b$ leave the same remainder after division by $m$. For example, $17\equiv5\pmod 4$, since both leave remainder $1$, and $-3\equiv1\pmod4$. The
remainder of $a$ on division by $m$ is called its \emph{residue} modulo $m$; it is one of $0,1,\dots,m-1$. Saying that
$a$ is odd is the same as saying $a\equiv1\pmod2$.

Residues can be added and multiplied consistently, as on a clock. The residues $0,1,\dots,m-1$ with this ``clock
arithmetic'' form a number system written $\Z/m\Z$. When $m=p$ is a prime we also write $\F_p$. In $\F_2=\{0,1\}$ we
have $1+1=0$, so adding and subtracting are the same thing.

For power series with integer coefficients, $U\equiv V\pmod m$ means that \emph{every} coefficient of $U-V$ is
divisible by $m$. Equivalently, the two series become equal when every coefficient is replaced by its residue; we call
this \emph{reducing} the series modulo $m$. The key fact is simple:

\emph{Reduction modulo $m$ is compatible with all the operations above.} To reduce a sum, a product, a substitution, or
an inverse, we may reduce the ingredients first and then compute with clock arithmetic. This holds because every
coefficient of the result is obtained from finitely many coefficients of the ingredients by adding, subtracting and
multiplying, and these operations are compatible with remainders.

\para{A word on ``rings''.} Several lemmas below hold not only for integers, but in every number system in
which one can add, subtract and multiply with the usual rules (such as $\Z$, $\Q$, and $\Z/m\Z$). Mathematicians call
such a number system a \emph{commutative ring}. We use the word only in this sense. No division is ever needed, which
matters because in $\Z/m\Z$ one cannot always divide (for instance, $2$ has no inverse in $\Z/4\Z$).

\subsection{The method in one paragraph}
Almost every proof in this paper follows the same three steps.
\begin{enumerate}
\item \emph{Uniqueness.} Show that the defining equation determines the coefficients one at a time, using only
addition, subtraction and multiplication. Then the equation has exactly one solution of the required shape, in every
commutative ring, including the clock arithmetic $\Z/2r\Z$.
\item \emph{Reduction.} Reducing the integer solution modulo $2r$ gives a solution of the reduced equation,
because reduction is compatible with all operations.
\item \emph{A good guess.} Find a simple series that satisfies the reduced equation and check that it does. By uniqueness,
this simple series \emph{is} the reduction of the sequence, and its coefficients can simply be read off.
\end{enumerate}
The creative part is step~3: finding a simple series that works. Step~1 sometimes needs care, when the equation seems to
require a division.
\subsection{A one-term Taylor formula}
\begin{lemma}\label{lem:taylor}
Let $\varepsilon$ be an element of a commutative ring with $\varepsilon^2=0$. If $u$ and $v$ are series with constant
term $0$, then for every series $S$
\[
S(u+\varepsilon v)=S(u)+\varepsilon\,v\,S'(u),
\]
where $S'=\sum_{n\ge1}n\,s_nx^{n-1}$ is the formal derivative of $S$.
\end{lemma}

\begin{proof}
By the binomial theorem, $(u+\varepsilon v)^j=u^j+j\,\varepsilon v\,u^{j-1}$, because all terms with
$\varepsilon^2$ or higher powers vanish. Multiply by $s_j$ and add up over $j$.
\end{proof}
An example of such an $\varepsilon$: in $\Z/8\Z$, the element $4$ satisfies $4^2=16=0$.


\section{Results}\label{sec:results}

For every integer $r\ge1$, Lemma~\ref{lem:unique} below shows that there is exactly one normalized series $F_r$ with
\begin{equation}\label{eq:family}
F_r=x+F_r^{\circ r}\cdot F_r^{\circ(r+1)},
\end{equation}
and that it has integer coefficients. The series of the entry is $A=F_{2}$. To state the result we use
two simple series with integer coefficients. For an integer $k$, let
\[
T_k(x)=\frac{x}{1-kx}=x+kx^2+k^2x^3+\cdots=\sum_{n\ge1}k^{\,n-1}x^n ,
\]
with the convention $k^0=1$ (so $T_0=x$), and let
\[
H(x)=\frac{x^4}{(1-x)^3}=x^4+3x^5+6x^6+10x^7+\cdots=\sum_{n\ge4}h_n\,x^n,\qquad h_n=\binom{n-2}{2}.
\]
We also set $h_n=0$ for $n<4$. (Lemma~\ref{lem:mobius} checks the formula for $h_n$.)

\begin{theorem}\label{thm:general}
Let $r\ge1$ and let $k$ be any integer. Modulo $2r$,
\[
F_r^{\circ k}\equiv
\begin{cases}
T_k+r\,H, & \text{if $r$ is even and $k$ is odd},\\
T_k, & \text{otherwise}.
\end{cases}
\]
In terms of coefficients: for every $n\ge1$,
\[
\coeff{n}{F_r^{\circ k}}\equiv k^{\,n-1}+r\,h_n\,[\,r\text{ even and }k\text{ odd}\,]\pmod{2r}.
\]
Moreover, $F_r^{\circ k}\equiv x\pmod{2r}$ if and only if $k$ is a multiple of $2r$. So, modulo $2r$, the series $F_r$
has order exactly $2r$ under composition.
\end{theorem}

In words: when $r$ is odd, the $k$-th iterate of $F_r$ has, modulo $2r$, the same coefficients as $x/(1-kx)$,
namely the powers $1,k,k^2,k^3,\dots$. When $r$ is even, the same is true for even $k$, while for odd $k$ the
coefficients are shifted by $r\,h_n$; since $2r\equiv0$, this shift only depends on whether $h_n$ is even or odd.

For the entry, the theorem gives the following.

\begin{corollary}\label{cor:app}
For $A=F_2$, the series of OEIS A396099, modulo $4$:
\begin{enumerate}[label=(\alph*)]
\item $a(n)$ is odd for every $n\ge1$;
\item for $n\ge3$, $a(n)\equiv1$ if $n\equiv2$ or $3\pmod4$, and $a(n)\equiv3$ if $n\equiv0$ or $1\pmod4$; so the
remainders run $1,3,3,1,1,3,3,1,\dots$ from $n=3$ on;
\item $A^{\circ2}\equiv x+2x^2\pmod4$, so $\coeff{n}{A(A(x))}\equiv0\pmod4$ for every $n\ge3$;
\item $A-x\,A^{\circ3}\equiv x+\dfrac{2x^3}{1-x}\pmod4$, so $\coeff{n}{\bigl(A(x)-x\,A(A(A(x)))\bigr)}\equiv2\pmod4$
for every $n\ge3$.
\end{enumerate}
\end{corollary}
These are exactly the four conjectures of the entry. Because the equation has only one normalized solution
(Lemma~\ref{lem:unique}), the four statements hold for every normalized series with integer coefficients that
satisfies \eqref{eq:def}; the Lean formalization states them in that form as well.

\section{Proofs}\label{sec:proofs}

\para{The idea in brief.} We follow the method of Section~\ref{sec:background}. First we show that
\eqref{eq:family} has exactly one solution, over every commutative ring, so it is enough to find \emph{some} solution
modulo $2r$. For odd $r$, the fraction $T_1=x/(1-x)$ works: its iterates are $T_k=x/(1-kx)$, and a one-line
calculation shows that $T_1$ satisfies the equation modulo $2r$. For even $r$, $T_1$ is not quite right, and we
correct it to $T_1+rH$. Because $r^2$ is a multiple of $2r$, every calculation with the correction can be done
``to first order'' with the Taylor formula of Lemma~\ref{lem:taylor}, which is exact here. The statements of the entry
are then read off from the formula.

\subsection{Existence and uniqueness}

\begin{lemma}\label{lem:unique}
Let $R$ be a commutative ring and $r\ge1$. There is exactly one normalized series $F$ with coefficients in $R$ such
that $F=x+F^{\circ r}\cdot F^{\circ(r+1)}$. Its coefficients are obtained from those of lower degree using only
additions and multiplications. In particular \eqref{eq:family} defines a series $F_r$ with integer coefficients, and
for every $m\ge2$ the reduction of $F_r$ modulo $m$ is the unique solution over $\Z/m\Z$.
\end{lemma}

\begin{proof}
Write $\Phi(F)=x+F^{\circ r}\cdot F^{\circ(r+1)}$ for the right-hand side. The key observation is:

\emph{If two normalized series $U$ and $V$ agree below degree $N$ (that is, $U-V=O(x^N)$, with $N\ge2$), then $\Phi(U)$
and $\Phi(V)$ agree below degree $N+1$.}

To see this, we first show that $U^{\circ k}-V^{\circ k}=O(x^N)$ for every $k\ge0$, by induction on $k$. The case
$k=0$ is clear. For the step, write
\[
U^{\circ(k+1)}-V^{\circ(k+1)}=\bigl(U(U^{\circ k})-U(V^{\circ k})\bigr)+\bigl(U(V^{\circ k})-V(V^{\circ k})\bigr).
\]
The first bracket is a sum of terms $u_j\bigl((U^{\circ k})^j-(V^{\circ k})^j\bigr)$, and each of them is divisible by
$U^{\circ k}-V^{\circ k}=O(x^N)$. The second bracket is the series $U-V=O(x^N)$ with $V^{\circ k}=x+\cdots$ substituted
for $x$, so it is $O(x^N)$ as well. Next, since all iterates start with $x$,
\[
U^{\circ r}U^{\circ(r+1)}-V^{\circ r}V^{\circ(r+1)}
=\bigl(U^{\circ r}-V^{\circ r}\bigr)\,U^{\circ(r+1)}+V^{\circ r}\,\bigl(U^{\circ(r+1)}-V^{\circ(r+1)}\bigr),
\]
and each product is $O(x^N)$ times a series starting with $x$, hence $O(x^{N+1})$. This proves the observation.

\emph{Existence.} Start with $G_0=x$ and put $G_{i+1}=\Phi(G_i)$. Each $G_i$ is normalized. The series $G_0=x$ and
$G_1=x+x^2+\cdots$ agree below degree $2$, and by the observation and induction, $G_i$ and $G_{i+1}$ agree below degree
$i+2$. So the coefficient of $x^n$ in $G_i$ no longer changes once $i\ge n-1$, and we let $F$ be the series whose
coefficient of $x^n$ is this final value. Then $F$ and $G_i$ agree below degree $i+2$, so $\Phi(F)$ and
$\Phi(G_i)=G_{i+1}$ agree below degree $i+3$; hence $\Phi(F)$ and $F$ agree below degree $i+2$ for every $i$, that is,
$\Phi(F)=F$.

\emph{Uniqueness.} If $U$ and $V$ are two normalized solutions, they agree below degree $2$. If they agree below
degree $N$, then $U=\Phi(U)$ and $V=\Phi(V)$ agree below degree $N+1$. So they agree in every degree.

The construction uses only additions and multiplications of coefficients, so it works over every commutative ring
and gives integer coefficients over $\Z$. Finally, reducing \eqref{eq:family} modulo $m$ shows that the reduction of
$F_r$ is a normalized solution over $\Z/m\Z$, because reduction is compatible with products and substitutions; by
uniqueness it is the only one.
\end{proof}

\subsection{Two calculations with fractions}

\begin{lemma}\label{lem:mobius}
Over every commutative ring, $T_c\circ T_d=T_{c+d}$ for all $c,d$. Consequently $T_1^{\circ k}=T_k$ for every integer
$k$, $T_0=x$, and $T_{-k}$ is the compositional inverse of $T_k$. Moreover $\coeff{n}{H}=\binom{n-2}{2}$ for $n\ge4$.
\end{lemma}

\begin{proof}
Substituting $T_d$ into $T_c$ and multiplying the numerator and the denominator by $1-dx$,
\[
T_c\bigl(T_d(x)\bigr)=\frac{x/(1-dx)}{1-c\,x/(1-dx)}=\frac{x}{1-dx-cx}=T_{c+d}(x).
\]
(Each step is an identity of power series, since all denominators have constant term $1$.) The other statements
follow: $T_1^{\circ k}=T_{1+\cdots+1}=T_k$ for $k\ge0$, and $T_k\circ T_{-k}=T_0=x$. For $H$: the coefficient of $x^j$
in $(1-x)^{-3}=(1+x+x^2+\cdots)^3$ is the number of ways to write $j$ as a sum of three nonnegative integers in order,
which is $\binom{j+2}{2}$; put $j=n-4$.
\end{proof}

The map $c\mapsto T_c$ thus turns addition of numbers into composition of series. This is why the iterates of $T_1$
are so simple. (The fractions $T_c$ are special cases of M\"obius transformations, whose compositions correspond to
products of $2\times2$ matrices.)

\begin{lemma}\label{lem:ident}
Let $T=T_1=x/(1-x)$ and let $T'=1/(1-x)^2$ be its derivative. Then, as power series with integer coefficients,
\begin{align}
H\cdot T'(T)+H(T)&=\frac{2x^4}{(1-2x)^3},\label{eq:id1}\\
(1-x)H-x^2\,T-x^3\,T'&=-\frac{2x^3}{1-x}.\label{eq:id2}
\end{align}
\end{lemma}

\begin{proof}
Since $1-T=\frac{1-2x}{1-x}$, we get $T'(T)=\frac{1}{(1-T)^2}=\frac{(1-x)^2}{(1-2x)^2}$ and
$H(T)=\frac{T^4}{(1-T)^3}=\frac{x^4}{(1-x)(1-2x)^3}$. Also $H\cdot T'(T)=\frac{x^4}{(1-x)(1-2x)^2}$. Adding,
\[
H\cdot T'(T)+H(T)=\frac{x^4\bigl((1-2x)+1\bigr)}{(1-x)(1-2x)^3}=\frac{2x^4}{(1-2x)^3}.
\]
For \eqref{eq:id2}, bring the three terms to the common denominator $(1-x)^2$: the numerator is
$x^4-x^3(1-x)-x^3=2x^4-2x^3=-2x^3(1-x)$.
\end{proof}

\subsection{Proof of Theorem~\ref{thm:general}}

We first prove the formula for $k\ge0$, separately for odd and even $r$, then extend it to negative $k$, and finally
prove the statement about the order.

\begin{proof}[Proof for odd $r$ and $k\ge0$]
Work modulo $2r$. We have
\[
(1-rx)\bigl(1-(r+1)x\bigr)=1-(2r+1)x+r(r+1)x^2\equiv1-x\pmod{2r},
\]
because $2r+1\equiv1$, and $r(r+1)=2r\cdot\frac{r+1}{2}$ is a multiple of $2r$ (here we use that $r+1$ is even). By
Lemma~\ref{lem:mobius}, $T^{\circ r}=T_r$ and $T^{\circ(r+1)}=T_{r+1}$, so
\[
x+T^{\circ r}\cdot T^{\circ(r+1)}=x+\frac{x^2}{(1-rx)\bigl(1-(r+1)x\bigr)}\equiv x+\frac{x^2}{1-x}=\frac{x}{1-x}=T\pmod{2r}.
\]
So $T$ satisfies the reduced equation, and by Lemma~\ref{lem:unique}, $F_r\equiv T\pmod{2r}$. Since reduction is
compatible with substitution, $F_r^{\circ k}\equiv T^{\circ k}=T_k$ for every $k\ge0$.
\end{proof}

\begin{proof}[Proof for even $r$ and $k\ge0$]
Work in $\Z/2r\Z$. There $r^2=2r\cdot\frac r2=0$, so Lemma~\ref{lem:taylor} applies with $\varepsilon=r$. Let
$f=T+rH$.

\emph{Step 1: $f^{\circ2}=T_2$.} By Lemma~\ref{lem:taylor},
\[
f(f)=T(T+rH)+r\,H(T+rH)=\bigl(T(T)+r\,H\,T'(T)\bigr)+r\,H(T)=T_2+r\bigl(H\,T'(T)+H(T)\bigr).
\]
(In the second term, Lemma~\ref{lem:taylor} would add $r\cdot r\,H\,H'(T)$, which is $0$ because $r^2=0$.) By
\eqref{eq:id1} the bracket equals $2x^4/(1-2x)^3$, and $r\cdot2=0$. So $f^{\circ2}=T_2$.

\emph{Step 2: $f$ satisfies the equation.} Since $r$ is even, Step~1 and Lemma~\ref{lem:mobius} give
$f^{\circ r}=(f^{\circ2})^{\circ r/2}=T_2^{\circ r/2}=T_r=x+rx^2+r^2x^3+\cdots=x+rx^2$. By Lemma~\ref{lem:taylor} again
(with $u=x$, $v=x^2$), $f^{\circ(r+1)}=f(x+rx^2)=f+rx^2f'$. Multiplying out and dropping every term that contains
$r^2$,
\[
f-x-f^{\circ r}\cdot f^{\circ(r+1)}=f-x-(x+rx^2)(f+rx^2f')=(1-x)f-x-rx^2f-rx^3f'.
\]
Now $(1-x)f=(1-x)T+r(1-x)H=x+r(1-x)H$, and in the terms already multiplied by $r$ we may replace $f$ by $T$ (the
difference is again a multiple of $r^2$). So the right side equals
\[
r\bigl((1-x)H-x^2T-x^3T'\bigr)=-\frac{2r\,x^3}{1-x}=0,
\]
by \eqref{eq:id2} and $2r=0$. Thus $f$ satisfies the reduced equation, and $F_r\equiv f\pmod{2r}$ by
Lemma~\ref{lem:unique}.

\emph{Step 3: all iterates with $k\ge0$.} If $k=2j$ is even, $f^{\circ k}=(f^{\circ 2})^{\circ j}=T_2^{\circ j}=T_k$.
If $k=2j+1$ is odd, $f^{\circ k}=f\circ T_{2j}=T(T_{2j})+r\,H(T_{2j})=T_k+r\,H(T_{2j})$. Over the integers,
$T_{2j}=x+2jx^2+\cdots$ agrees with $x$ modulo $2$, so $H(T_{2j})-H$ has even coefficients, and $r$ times it is a
multiple of $2r$. Hence $r\,H(T_{2j})=r\,H$ in $\Z/2r\Z$, and $f^{\circ k}=T_k+rH$.
\end{proof}

\begin{proof}[Proof for negative $k$, and of the statement about the order]
Write $E_k$ for the series on the right side of the theorem ($T_k+rH$ or $T_k$). For $k=2r$ we have shown
$F_r^{\circ 2r}\equiv E_{2r}=T_{2r}$, and $T_{2r}=x+2r\,x^2+(2r)^2x^3+\cdots\equiv x$. So $F_r^{\circ2r}\equiv x$, and
therefore also its inverse $F_r^{\circ(-2r)}\equiv x$. Now let $k$ be any integer, and write $k=j+2rq$ with
$0\le j<2r$ and $q$ an integer. Then $F_r^{\circ k}=F_r^{\circ j}\circ(F_r^{\circ 2r})^{\circ q}\equiv F_r^{\circ j}\equiv E_j$.
Finally $E_j\equiv E_k$: the coefficients $j^{n-1}$ and $k^{n-1}$ of $T_j$ and $T_k$ are congruent modulo $2r$
because $j\equiv k$, and $j$ and $k$ have the same parity because $2r$ is even. This proves the formula for every
integer $k$. The coefficient form follows from $\coeff{n}{T_k}=k^{n-1}$ and Lemma~\ref{lem:mobius}.

For the order: the coefficient of $x^2$ in $E_k$ is $k$, because $H$ starts with $x^4$. So if $F_r^{\circ k}\equiv x$,
then $k\equiv0\pmod{2r}$. Conversely, if $2r$ divides $k$, then $k$ is even, so $E_k=T_k$, and every coefficient
$k^{n-1}$ with $n\ge2$ is a multiple of $k$, hence of $2r$; so $F_r^{\circ k}\equiv T_k\equiv x$.
\end{proof}

\subsection{Proof of Corollary~\ref{cor:app}}

\begin{proof}[Proof of Corollary~\ref{cor:app}]
Theorem~\ref{thm:general} with $r=2$ gives $A\equiv T_1+2H$ and $A^{\circ2}\equiv T_2\pmod4$.

(a) and (b). The coefficients of $T_1$ are all $1$, so $a(n)\equiv1+2h_n\pmod4$. This is odd, which proves (a). It
is $1$ when $h_n$ is even and $3$ when $h_n$ is odd. The parity of $h_n=\binom{n-2}{2}$ repeats with period $4$,
because $h_{n+4}-h_n=\binom{n+2}{2}-\binom{n-2}{2}=4n-2$ is even, and $h_2,h_3,h_4,h_5=0,0,1,3$. So $h_n$ is odd
exactly when $n\equiv0$ or $1\pmod4$, which proves (b).

(c). $T_2=x+2x^2+4x^3+8x^4+\cdots\equiv x+2x^2\pmod4$.

(d). Work modulo $4$ and let $f=T+2H$ be the reduction of $A$. By Lemma~\ref{lem:taylor} with $\varepsilon=2$
(note $2^2=0$ modulo $4$), $f^{\circ3}=f(f^{\circ2})=f(x+2x^2)=f+2x^2f'$. Hence, using $(1-x)T=x$ and dropping
multiples of $4$,
\[
f-x\,f^{\circ3}=(1-x)f-2x^3f'=x+2\bigl((1-x)H-x^3T'\bigr)=x+2\,\frac{x^4-x^3}{(1-x)^2}=x-\frac{2x^3}{1-x}.
\]
Since $-2\equiv2\pmod4$, this is $x+2x^3/(1-x)=x+2x^3+2x^4+\cdots$.
\end{proof}

\section{Checking the results}\label{sec:checks}

A mathematical proof can contain mistakes, so the results were checked in two independent ways.

\para{Numerical check.}
We computed the coefficients of $A$ up to $x^{23}$ exactly from \eqref{eq:def}. They reproduce all
$21$ terms published in the entry. Working modulo $8$, we then computed $A$, $A^{\circ2}$ and $A^{\circ3}$ up to
$x^{159}$ and confirmed statements (a)--(d) of Corollary~\ref{cor:app} for every $n<160$. Finally, we confirmed the
formula of Theorem~\ref{thm:general}, modulo $2r$, for every $r\le8$, every $k$ with $-4\le k\le9$, and every $n<90$.
The computation was done with the PARI/GP computer algebra system, using code that was written separately from the
formal proofs (file \texttt{verification/independent\_checks.gp} in the accompanying repository~\cite{repo}).
A computation can only test finitely many cases, so it is a check on the definitions and on the indexing, not a proof.

\para{Formal proof in Lean.}
Lean~\cite{lean} is a \emph{proof assistant}: a program in which mathematical statements and their proofs are written
in a precise formal language, and every step is checked mechanically down to the basic rules of logic. All results of
this paper are proved in Lean~4 (version 4.19.0) on top of Mathlib~\cite{mathlib}, Lean's library of formalized
mathematics. The statements are collected in the file \texttt{lean/Entries/A396099.lean} of the repository,
in the namespace \texttt{FinalCollection.A396099}. The general results for the whole family are proved in the folder \texttt{lean/Consecutive} of the
repository; the last rows of the table give their full names. In Lean, iterates with negative index are built from the
compositional inverse, exactly as in Section~\ref{sec:background}.
The table lists where each result of this paper is proved.
\begin{center}\small\renewcommand{\arraystretch}{1.25}
\begin{tabular}{>{\raggedright\arraybackslash}p{0.52\linewidth}>{\raggedright\arraybackslash}p{0.4\linewidth}}
\toprule
Statement in this paper & Lean theorem(s)\\
\midrule
Lemma~\ref{lem:unique} for $r=2$; the series $A$ is normalized and satisfies \eqref{eq:def} & \texttt{source\_\allowbreak{}existence\_\allowbreak{}uniqueness}, \texttt{source\_\allowbreak{}normalized}, \texttt{source\_\allowbreak{}equation} \\
The order: $A^{\circ k}\equiv x\pmod4$ if and only if $4$ divides $k$ ($k\in\Z$) & \texttt{exact\_\allowbreak{}composition\_\allowbreak{}order} \\
Corollary~\ref{cor:app}(a)--(d) for every normalized integer solution of \eqref{eq:def} & \texttt{a396099\_\allowbreak{}all\_\allowbreak{}four} \\
Corollary~\ref{cor:app}(a) and (b) & \texttt{a396099\_\allowbreak{}odd}, \texttt{a396099\_\allowbreak{}residue}, \texttt{a396099\_\allowbreak{}period} \\
Corollary~\ref{cor:app}(c): $A^{\circ2}\equiv x+2x^2$, and its coefficients & \texttt{a396099\_\allowbreak{}second\_\allowbreak{}iterate\_\allowbreak{}series}, \texttt{a396099\_\allowbreak{}second\_\allowbreak{}iterate} \\
Corollary~\ref{cor:app}(d): $A-xA^{\circ3}\equiv x+2x^3/(1-x)$, and its coefficients & \texttt{a396099\_\allowbreak{}difference\_\allowbreak{}series}, \texttt{a396099\_\allowbreak{}difference} \\
Lemma~\ref{lem:unique} for every $r\ge1$, over every commutative ring & \texttt{Consecutive.existence\_\allowbreak{}uniqueness} \\
Theorem~\ref{thm:general}, the formula for every $r\ge1$ and every integer $k$ & \texttt{Consecutive.ziter\_\allowbreak{}solution\_\allowbreak{}mod} \\
Theorem~\ref{thm:general}, the statement about the order, for every $r\ge1$ & \texttt{Consecutive.ziter\_\allowbreak{}solution\_\allowbreak{}mod\_\allowbreak{}eq\_\allowbreak{}X\_\allowbreak{}iff} \\
\bottomrule
\end{tabular}
\end{center}
The Lean project was rebuilt from source for this paper, and Lean was asked to list the basic assumptions (axioms)
that each of these theorems ultimately depends on. Every one depends only on the three standard axioms of Lean's logic
(\texttt{propext}, \texttt{Classical.choice}, \texttt{Quot.sound}); no step is left unproved. A formal proof certifies
exactly the statement that Lean checks. Whether that statement says what the OEIS entry means is a separate question,
so the Lean definitions were also compared, side by side, with the definition in the OEIS entry
(Section~\ref{sec:sequence}).

\section{Discussion}\label{sec:discussion}

\para{Scope.} Theorem~\ref{thm:general} determines every coefficient of every iterate of $A$ modulo $4$, so it
settles all four conjectures of the entry completely. Theorem~\ref{thm:general} says nothing about remainders modulo numbers that do
not divide $2r$, or about the size of the coefficients.

\para{All coefficients are odd.} Since $2$ divides $2r$, the theorem also determines every $F_r$ modulo $2$: the
correction $rH$ has even coefficients when $r$ is even, so in all cases $F_r\equiv x/(1-x)\pmod2$. Thus every
coefficient of every $F_r$ is odd.

\para{Higher moduli.} A computer calculation suggests that the remainders modulo $4r$ are periodic too. For $r=1,2,3,4,5$
and all $n<130$, the residues of $\coeff{n}{F_r}$ modulo $4r$ repeat with period $3$, $4$, $6$, $8$ and $3$
respectively (from $n=1$, $4$, $1$, $2$ and $1$ on). We have not tried to prove this, and we do not know a formula for
the period in general.

\para{A more general equation.} The proof for odd $r$ uses very little. Consider, for any integers $r,s\ge1$, the
equation $F=x+F^{\circ r}\cdot F^{\circ s}$. The proof of Lemma~\ref{lem:unique} works without change, so it has exactly
one normalized solution over every commutative ring. Let $m$ be a common divisor of $r+s-1$ and $rs$. Then
$(1-rx)(1-sx)=1-(r+s)x+rs\,x^2\equiv1-x\pmod m$, and the calculation in the proof for odd $r$ shows that
$x/(1-x)$ solves the equation modulo $m$. Hence every iterate of the solution is congruent to $x/(1-kx)$ modulo $m$.
For $s=r+1$ with $r$ odd this is the case $m=2r$ of Theorem~\ref{thm:general}; another example is $r=1$, where
$F=x+F\cdot F^{\circ s}$ gives $F\equiv x/(1-x)\pmod s$. We checked this remark by computer for all $r\le7$ and
$s\le8$ with $m>1$ ($38$ pairs), for $-3\le k\le6$ and all $n<50$. It is not part of the Lean formalization.

\para{Related work.} Normalized power series with composition as the operation form a \emph{group}: composition is
associative, $x$ is a neutral element, and every element has an inverse (Lemma~\ref{lem:inverse}). When the
coefficients are taken modulo a prime $p$, this group is known as the \emph{Nottingham group}. It has been studied
extensively in group theory, and in particular its elements of finite order have been investigated in
detail~\cite{camina}. Theorem~\ref{thm:general} produces, from a natural equation, an explicit series of order exactly
$2r$ in the corresponding group with coefficients modulo $2r$.

\section*{About this paper}\label{sec:about}
\para{How the work was done.}
This paper is one of 34 papers on conjectures from the OEIS that came out of an autonomous AI research effort directed
by Dev~D.~Goyal~\cite{repo}. The roles were as follows.
\begin{itemize}
\item \emph{Dev~D.~Goyal} started the project, chose which problems to pursue, asked for proofs, formal verification and
revisions, and decided what to publish. Dev~D.~Goyal did not supply proof strategies or intermediate arguments.
\item \emph{ChatGPT} (OpenAI), using the GPT-6 model with Pro reasoning, worked autonomously from those prompts. It
searched the OEIS for suitable conjectures, found the proofs, wrote the Lean~4 formalization, and wrote a first draft
of this paper.
\item \emph{Claude Opus~5.5}, used with Claude Code, checked the work independently and then rewrote this paper. It
rebuilt the Lean project from source and confirmed that the theorems listed in Section~\ref{sec:checks} rely only on
Lean's standard axioms; compared the Lean definitions with the OEIS definition; recomputed the sequence with separately
written PARI/GP code, compared it with the terms published in the OEIS, and tested every claim of the paper on many
terms; and read every step of the proofs. It then rewrote the text to make it self-contained and readable. The
theorems and the proof strategies are ChatGPT's. In rewriting, Claude Opus~5.5 expanded the explanations and, in
places, wrote out an argument in more detail or by a more elementary route. The Discussion section was written during
this rewrite; the further observations it contains were checked by computer, and they are not part of the Lean
formalization.
\end{itemize}
Dev~D.~Goyal has verified the results and is responsible for this paper.

\para{Review status.} This paper has not been peer reviewed, and it has not been submitted to a journal or to the OEIS.

\para{Priority.} An OEIS entry that labels a statement as a conjecture shows only that its contributor did not know
a proof; it does not show that no proof existed. During the research only limited searches for earlier work were
made. No claim of priority is made.

\para{Corrections.} Please report any error in the issue tracker of the repository~\cite{repo}, where the Lean
sources, the numerical checks, and the earlier drafts are also available.

\section*{Acknowledgements}
The sequence and its conjectures are due to Paul~D.~Hanna~\cite{oeis}, who also contributed the other members of the
family: A396797, A396798, A396807 (the four entries correspond to $r=2,3,4,5$). The case $r=3$ (entry A396797) was
proved earlier by Shiva~Priyan~Shanmuga~Sundaram~\cite{sundaram}. 

\begin{thebibliography}{99}\small
\raggedright
\bibitem{oeis} P.~D.~Hanna, entry A396099 (G.f.\ satisfies $A(x)=x+A^{2}(x)\cdot A^{3}(x)$), \emph{The On-Line Encyclopedia of Integer Sequences}, June 2026. \url{https://oeis.org/A396099}. Text as of October 2026.
\bibitem{sundaram} S.~P.~Shanmuga~Sundaram, \emph{A proof of the congruence conjectures for OEIS A396797}, September 2026. Linked from the OEIS entry: \url{https://oeis.org/A396797/a396797.pdf}.
\bibitem{camina} R.~Camina, The Nottingham group, in \emph{New Horizons in pro-$p$ Groups}, M.~du~Sautoy, D.~Segal and A.~Shalev (eds.), Progress in Mathematics 184, Birkh\"auser, 2000, 205--221.
\bibitem{oeisdb} N.~J.~A.~Sloane and The OEIS Foundation Inc., \emph{The On-Line Encyclopedia of Integer Sequences}. \url{https://oeis.org}.
\bibitem{stanley} R.~P.~Stanley, \emph{Enumerative Combinatorics}, Volume~2, Cambridge University Press, 1999. (Chapter~5 treats formal power series, composition and compositional inverses.)
\bibitem{niven} I.~Niven, Formal power series, \emph{American Mathematical Monthly} 76 (1969), 871--889.
\bibitem{lean} L.~de~Moura and S.~Ullrich, The Lean~4 theorem prover and programming language, in \emph{Automated Deduction -- CADE 28}, Lecture Notes in Computer Science 12699, Springer, 2021, 625--635.
\bibitem{mathlib} The mathlib Community, The Lean mathematical library, in \emph{Proceedings of the 9th ACM SIGPLAN International Conference on Certified Programs and Proofs (CPP 2020)}, 367--381. Version used: commit \texttt{c44e0c8ee63c}.
\bibitem{repo} D.~D.~Goyal, \emph{Proofs and corrections for 34 OEIS entries, with Lean~4 formalizations}, repository, 2026. \url{https://github.com/dvgyl/oeis-a396843-proof}.
\end{thebibliography}
\end{document}
